\documentclass{article}
\begin{document}
\section{はじめに}
従来の目視検査はコストが高く、効率が低い。自動化された病害検出手法が強く求められている。
\section{検出手法の設計}
本手法は二段階アーキテクチャを採用する。第一段階で画像の前処理と強化を行い、第二段階で病害分類を実行する。
粒子群最適化により各パラメータを自動探索する。粒子群最適化は大域探索に優れ、局所解に陥りにくい。
データ拡張戦略にはランダムクロップ、色ジッター、弾性変形が含まれる。
\section{実験と結果分析}
本手法の精度は91.3%に達し、ベースラインを2.4ポイント上回った。訓練には50エポックを要した。
\end{document}
